\clearpage
For claims 6--10, wealth lotteries have finite positive outcomes. Bernoulli
indices are strictly increasing; use risk-aversion coefficients where defined.
\task{6. A change of utility units}{}
Replacing $b(w)$ by $3b(w)+7$ preserves expected-utility preferences,
certainty equivalents and Arrow--Pratt absolute risk aversion where defined.
\task{7. Increasing is enough}{}
If $b(w)>0$, replacing it by $[b(w)]^2$ preserves the ranking of all wealth
lotteries, because squaring is increasing on positive numbers.
\task{8. Dominance without concavity}{}
If lottery A first-order stochastically dominates lottery B, every increasing
Bernoulli index weakly prefers A; concavity is unnecessary.
\task{9. Decreasing risk aversion prevents full insurance}{}
A consumer with decreasing absolute risk aversion cannot optimally choose
full coverage at an actuarially fair insurance premium.
\task{10. One zero premium reveals everything}{}
For a strictly increasing concave Bernoulli index, a zero risk premium for one
nondegenerate lottery implies risk neutrality over all positive-wealth lotteries.
